\section{Recap}
\begin{frame}{}
    \LARGE Policy Optimization: \textbf{Recap (Day 3)}
\end{frame}

\begin{frame}{Recap: What Day 3 Established}
\begin{itemize}
    \item \textbf{The objective is unchanged from Day 3} --- maximize the expected discounted return:
    $$\mathcal{J}(\theta) = \mathbb{E} \left [ \sum_{t \geq 0} \gamma^t r_t \mid \pi_\theta \right ], \qquad \theta^{\star} = \arg\max_\theta \mathcal{J}(\theta)$$
    \item Day 3 derived its gradient and, after variance reduction, ended at the estimator:
    $$\nabla_\theta \mathcal{J}(\theta) \approx \sum_{t \geq 0} A^{\pi}(s_t, a_t)\, \nabla_\theta \log \pi_\theta(a_t|s_t)$$
    \item $A^{\pi}(s,a) = Q^{\pi}(s,a) - V^{\pi}(s)$ is the \textbf{advantage function} --- the best baseline
    \item Using $A^{\pi}$ instead of raw returns reduces variance without introducing bias
    \item \textbf{Unresolved question (Day 3's closing slide):} how do we actually estimate $A^{\pi}$?
\end{itemize}
\end{frame}

\begin{frame}{Recap: What Today Delivers}
\begin{itemize}
    \item \textbf{Today we cash that promise.}
    \item We learn a critic $V_\phi$ alongside the policy --- this is the \textbf{Actor-Critic} architecture
    \item Good news: we only need to learn $V^{\pi}$, not $Q^{\pi}$ --- one network, one output per state
    \item The critic is trained by \textbf{TD learning} (which we unpack in a few slides); its one-step error is what gives us an online, low-variance estimate of $A^{\pi}$
    \item Then we scale up: A2C/A3C $\rightarrow$ trust region (TRPO) $\rightarrow$ clipped objective (PPO)
\end{itemize}
\end{frame}
